\section{Results}
\label{sec:results}

\subsection{Primary Ablation Results}

\Cref{tab:main} reports \valbpb{} improvement for each group across three
independent repeats.

\begin{table}[!htbp]
\centering
\caption{%
  \valbpb{} change ($\Delta = \text{best} - \text{baseline}$, more negative
  is better) over 30 inner iterations.
  Baseline \valbpb{} varies slightly across repeats due to training
  randomness (range 1.094--1.114).
  Group~C's mean improvement is 5$\times$ that of Group~A and 7.5$\times$ that
  of Group~B (by absolute $|\Delta|$).
  Bold values indicate the best repeat within Groups~C and~D.%
}
\label{tab:main}
\begin{tabular}{@{}lrrrr@{}}
\toprule
Group & R1 & R2 & R3 & Mean $\pm$ Std \\
\midrule
A (Level~1)         & $-0.009$ & $-0.008$ & $-0.011$ & $-0.009 \pm 0.002$ \\
B (Level~1 + 1.5)   & $-0.000$ & $-0.010$ & $-0.009$ & $-0.006 \pm 0.006$ \\
C (Level~1 + 1.5 + 2) & $\mathbf{-0.065}$ & $-0.011$ & $\mathbf{-0.058}$ & $\mathbf{-0.045 \pm 0.030}$ \\
D (Level~1 + 2) & $-0.001$ & $\mathbf{-0.063}$ & $-0.039$ & $-0.034 \pm 0.031$ \\
\bottomrule
\end{tabular}
\end{table}

Group~A achieves consistent but small improvements: $-0.009 \pm 0.002$.
Group~B is comparable to Group~A ($-0.006 \pm 0.006$); its R1 found
essentially no improvement ($-0.000$), inflating variance.
Group~C achieves $-0.045 \pm 0.030$, a 5$\times$ improvement over Group~A.
Two of three Group~C repeats (R1 and R3) show dramatic gains ($-0.065$ and
$-0.058$); R2 underperformed at $-0.011$.
Group~D achieves $-0.034 \pm 0.031$: R2 reached the best single result
across all D repeats ($-0.063$) while R1 barely improved ($-0.001$), giving
a mean comparable to, but slightly below, Group~C.
This confirms that Level~2 is the primary driver of improvement and that
Level~1.5 is not essential when Level~2 is present.
\cref{fig:convergence} shows the search trajectories across all twelve runs.

\begin{figure}[!htbp]
  \centering
  \includegraphics[width=\columnwidth]{convergence.pdf}
  \caption{%
    Running-minimum \valbpb{} vs.\ iteration for all 9 runs.
    Thin lines: individual repeats; thick lines: group means.
    Vertical dashed lines mark Level~2 interventions at iterations 10 and 20 (Group~C only).
    Group~C's R1 and R3 show a sharp drop after Level~2 mechanisms
    guide the search toward \texttt{TOTAL\_BATCH\_SIZE} reduction.%
  }
  \label{fig:convergence}
\end{figure}

\subsection{Level~2 Mechanism Inventory}

\Cref{tab:mechanisms} lists all mechanisms generated by Level~2 across the
six research sessions (two per repeat in Group~C).

\begin{table}[!htbp]
\centering
\caption{%
  Level~2 mechanism inventory.
  All code was generated on the first attempt (zero retries).
  Five of six mechanisms passed import validation and were activated;
  one (GP Regressor) was reverted due to a missing \texttt{sklearn}
  dependency.%
}
\label{tab:mechanisms}
\begin{tabular}{@{}llllcc@{}}
\toprule
Repeat & Round & Mechanism & Domain & Import & Active \\
\midrule
R1 & 1 & Tabu Search Manager     & Combinatorial opt.    & \checkmark & \checkmark \\
R1 & 2 & Helper class            & ---                   & \checkmark & \checkmark \\
R2 & 1 & Multi-Scale Bandit Proposer & Online learning / MAB & \checkmark & \checkmark \\
R2 & 2 & Helper class            & ---                   & \checkmark & \checkmark \\
R3 & 1 & GP Regressor            & Bayesian optimization & $\times^{*}$ & reverted \\
R3 & 2 & Syst.\ Orthogonal Exploration & DOE        & \checkmark & \checkmark \\
\bottomrule
\multicolumn{6}{@{}l}{%
  $^{*}$ Code valid but \texttt{sklearn} not installed; reverted automatically.%
}
\end{tabular}
\end{table}

The two ``Helper class'' entries are utility classes generated alongside the primary
mechanisms (e.g., data structures for tracking parameter history); they are auxiliary
to the named mechanism in each round.

The three active named mechanisms are:
\textbf{Tabu Search Manager} (maintains a tabu list of recently visited
parameter regions, preventing the LLM from reproposing the same changes);
\textbf{Multi-Scale Bandit Proposer} (treats parameter selection as a
multi-armed bandit, balancing exploration and exploitation across parameters
at different scales); and
\textbf{Systematic Orthogonal Exploration} (forces the LLM to explore
orthogonal parameter dimensions, preventing over-focus on a single parameter).
Each mechanism was drawn from a different domain; Level~2 was not told which
domains to consider.

\subsection{Search Behavior Analysis}

The three groups exhibit qualitatively different search trajectories.

\paragraph{Group~A: near-deterministic repetition.}
All three repeats follow nearly the same proposal sequence from the same
baseline:
iteration~1 attempts \texttt{TOTAL\_BATCH\_SIZE} increase (discard);
iteration~2 reduces \texttt{WEIGHT\_DECAY} (keep, $\Delta \approx -0.008$);
iteration~3 sets \texttt{WINDOW\_PATTERN="SSSS"} (keep, $\Delta \approx -0.002$);
iterations~4--30 repeat these same two changes, accumulating up to 22
consecutive discards.
The LLM, given the same state, generates nearly the same proposals every time.

\paragraph{Group~B: redirected but bounded.}
Level~1.5 correctly identifies stalled parameters and redirects search:
by cycle~2--3 the outer loop freezes \texttt{WEIGHT\_DECAY} and
\texttt{WINDOW\_PATTERN} and redirects toward \texttt{LR},
\texttt{UNEMBEDDING\_LR}, \texttt{MATRIX\_LR}, and \texttt{FINAL\_LR\_FRAC}.
Group~B explores more parameters than Group~A, but achieves similarly sized
improvements because it operates within the same structural keep/discard
framework.

\paragraph{Group~C: Level~2 mechanisms unlock new directions.}
The decisive event in Group~C's R1 and R3 is the discovery of
\texttt{TOTAL\_BATCH\_SIZE} \emph{reduction} (from $2^{19}$ to $2^{17}$ or
$2^{18}$), which produces improvements of $-0.039$ to $-0.065$---roughly
5--8$\times$ larger than any single change found by Groups~A or~B.

\paragraph{Group~D: Level~2 without outer loop guidance.}
Group~D's pattern is similar to Group~C: R2 and R3 independently discover
\texttt{TOTAL\_BATCH\_SIZE} reduction (D2 reaches $2^{17}$, D3 reaches $2^{18}$),
producing improvements of $-0.063$ and $-0.039$.
However, R1 failed to benefit: its two Level~2 sessions generated mechanisms
(\texttt{diversity\_enforcer} and \texttt{fixation\_detector}) that failed import
validation, leaving it to run as bare Level~1 with no mechanism injection; the
inner loop could not discover the batch size direction on its own.
The absence of Level~1.5 means there is no focused parameter-freeze guidance
to steer Level~2's attention, so mechanism quality is more variable across
repeats---explaining the higher variance ($\pm 0.031$) relative to Group~C's
repeats where Level~1.5 provides enriched trace context.

\subsection{The \texttt{TOTAL\_BATCH\_SIZE} Discovery}

The most impactful finding across all experiments is that reducing
\texttt{TOTAL\_BATCH\_SIZE} from $2^{19}$ to $2^{17}$--$2^{18}$ dramatically
improves \valbpb{} on the RTX~5090 under a 300-second training budget.
The mechanism is straightforward: a smaller batch size yields more gradient
steps within the fixed time budget, and better convergence for this 50M-parameter
model.
The original $2^{19}$ batch size was tuned for H100 throughput; the RTX~5090
running SDPA (Flash Attention~3 is unsupported on Blackwell compute~12.0) has
different optimal batch characteristics.

Groups~A and~B both miss this direction for the same reason: DeepSeek's default
search path attempts \texttt{TOTAL\_BATCH\_SIZE} \emph{increase} first (an
implicit ``larger batch is better'' bias).
After the increase is discarded, Group~A repeats it; Group~B's outer loop
freezes \texttt{TOTAL\_BATCH\_SIZE} after the failed increase, blocking the
decrease direction entirely.
Only Group~C's Level~2 mechanisms---specifically, Tabu Search (which prevents
revisiting failed directions) and Orthogonal Exploration (which forces
dimensional diversity)---pushed the LLM to try the decrease direction.
